\chapter{Disaster Management and Hazard Mitigation}
\label{chap:disaster-management}

\section{Overview and Course Learning Outcomes}
Disasters represent severe disruptions of community functioning where human, material, economic, and environmental losses exceed the coping capacity of affected populations. Disasters emerge from the intersection of natural or technological hazards with human exposure and physical-social vulnerability. This chapter examines the disaster risk formulation, geophysical and hydro-meteorological hazards, man-made industrial disasters, the National Disaster Management Act 2005, and the global Sendai Framework for Disaster Risk Reduction (2015--2030).

\begin{important}[Course Learning Outcomes for Unit 5]
Upon completing this chapter, the student will be capable of:
\begin{enumerate}[leftmargin=1.5em]
    \item Formulating and computing the Disaster Risk Equation ($R = [H \times V \times E] / C$).
    \item Comparing geophysical hazards (earthquakes, landslides, tsunamis) and calculating seismic energy via Gutenberg-Richter equations.
    \item Analyzing hydro-meteorological disasters (tropical cyclones, flash floods, droughts) and evaluating IMD warning scales.
    \item Diagnosing chemical and industrial technological disasters through the Bhopal Gas Tragedy case study.
    \item Outlining the 3-tier Indian disaster management architecture (NDMA, SDMA, DDMA) and the 4 Sendai Framework priorities.
\end{enumerate}
\end{important}

\section{The Disaster Risk Formulation}
\begin{definition}[Hazard, Vulnerability, and Disaster Risk]
\begin{itemize}[leftmargin=1.2em]
    \item \textbf{Hazard ($H$):} A dangerous physical phenomenon, substance, or human condition that may cause loss of life, injury, property damage, or environmental degradation.
    \item \textbf{Exposure ($E$):} The presence of people, infrastructure, housing, and production capacities in hazard-prone zones.
    \item \textbf{Vulnerability ($V$):} The structural, social, economic, or geographic conditions that increase the susceptibility of a community to the impacts of hazards.
    \item \textbf{Capacity ($C$):} The combination of all strengths, attributes, and resources available within an organization or community to manage and reduce disaster risks.
\end{itemize}
\end{definition}

\begin{keyequation}[Universal Disaster Risk Equation]
Disaster risk is quantified as a direct function of hazard, exposure, and vulnerability, inversely moderated by societal coping capacity:
\begin{equation}
\text{Disaster Risk} = \frac{\text{Hazard} \times \text{Exposure} \times \text{Vulnerability}}{\text{Capacity}}
\end{equation}
Disasters occur only when a physical hazard strikes an exposed, vulnerable population with insufficient coping mechanisms.
\end{keyequation}

\section{Geophysical and Atmospheric Disasters}
\subsection{Earthquakes and Seismic Energy Release}
Earthquakes arise from sudden tectonic shear failure along geological faults, releasing accumulated elastic strain energy as seismic waves ($P$-waves, $S$-waves, Love and Rayleigh surface waves).

\begin{keyequation}[Gutenberg-Richter Energy Release Equation]
The total radiated seismic energy $E$ (in Joules) is related to Richter local magnitude $M$ by:
\begin{equation}
\log_{10} E = 4.8 + 1.5 M \implies E = 10^{4.8 + 1.5 M}
\end{equation}
Consequently, a unit increase in magnitude ($\Delta M = 1$) corresponds to a $10^{1.5} = \sqrt{1000} \approx 31.62$-fold increase in radiated mechanical energy. A magnitude $7.0$ earthquake releases $1{,}000$ times more energy than a magnitude $5.0$ event.
\end{keyequation}

\subsection{Tropical Cyclones and IMD Classification}
Tropical cyclones are intense low-pressure vortex systems fueled by latent heat of condensation over warm tropical ocean waters (sea surface temperature $\ge 26.5^\circ\text{C}$). The India Meteorological Department (IMD) classifies cyclonic disturbances by maximum sustained 3-minute surface winds:
\begin{itemize}[leftmargin=1.5em]
    \item \textbf{Cyclonic Storm:} $62\text{--}88\text{ km/h}$ ($34\text{--}47\text{ knots}$).
    \item \textbf{Severe Cyclonic Storm:} $89\text{--}117\text{ km/h}$ ($48\text{--}63\text{ knots}$).
    \item \textbf{Very Severe Cyclonic Storm:} $118\text{--}166\text{ km/h}$ ($64\text{--}89\text{ knots}$).
    \item \textbf{Extremely Severe Cyclonic Storm:} $167\text{--}221\text{ km/h}$ ($90\text{--}119\text{ knots}$).
    \item \textbf{Super Cyclonic Storm:} $\ge 222\text{ km/h}$ ($\ge 120\text{ knots}$) (e.g. 1999 Odisha Super Cyclone).
\end{itemize}

\section{Technological and Chemical Disasters}
\begin{casestudy}[The Bhopal Gas Disaster (December 2--3, 1984)]
The world's worst industrial disaster occurred at the Union Carbide India Limited (UCIL) pesticide formulation plant in Bhopal, Madhya Pradesh. Over $42\text{ tonnes}$ of toxic Methyl Isocyanate (MIC, $CH_3N=C=O$) gas escaped from Tank E610 following runaway exothermic hydration when water entered the storage tank. Multiple passive and active safety barriers had been decommissioned or failed:
\begin{itemize}[leftmargin=1.2em]
    \item Refrigeration chiller unit was turned off to save operational costs.
    \item Caustic soda vent gas scrubber was non-operational.
    \item Flare tower piping was disconnected for maintenance.
    \item Water spray curtains had insufficient hydraulic pressure to reach the stack discharge height.
\end{itemize}
Thermal inversion trapped the dense, toxic MIC gas close to the ground, dispersing it across adjacent impoverished informal settlements (Chhola, Kenwood). Over $3{,}800$ people died immediately from acute pulmonary edema, chemical asphyxiation, and respiratory failure, with long-term chronic morbidity afflicting over $500{,}000$ survivors. This catastrophic failure prompted the enactment of the comprehensive \textbf{Environment (Protection) Act, 1986} and the \textbf{Public Liability Insurance Act, 1991}.
\end{casestudy}

\section{Disaster Management Framework: DM Act 2005 and Sendai Framework}
India transitioned from a post-disaster relief-centric model to a proactive, holistic prevention and mitigation paradigm through the enactment of the \textbf{Disaster Management Act, 2005}:
\begin{itemize}[leftmargin=1.5em]
    \item \textbf{National Level:} National Disaster Management Authority (NDMA), chaired ex-officio by the Prime Minister of India.
    \item \textbf{State Level:} State Disaster Management Authority (SDMA), chaired by the respective Chief Minister.
    \item \textbf{District Level:} District Disaster Management Authority (DDMA), headed by the District Magistrate / Collector.
    \item \textbf{Operational Response:} National Disaster Response Force (NDRF), comprising specialized battalions trained for search, rescue, and chemical-biological-radiological-nuclear (CBRN) emergencies.
\end{itemize}

\begin{lawbox}{The Sendai Framework for Disaster Risk Reduction (2015--2030)}
Adopted at the Third UN World Conference in Sendai, Japan, it establishes 4 Priorities for Action:
\begin{enumerate}[leftmargin=1.2em]
    \item \textbf{Priority 1:} Understanding disaster risk in all its dimensions of vulnerability and exposure.
    \item \textbf{Priority 2:} Strengthening disaster risk governance to manage disaster risk.
    \item \textbf{Priority 3:} Investing in disaster risk reduction for structural and non-structural resilience.
    \item \textbf{Priority 4:} Enhancing disaster preparedness for effective response and to \textbf{``Build Back Better''} in recovery, rehabilitation and reconstruction.
\end{enumerate}
\end{lawbox}

\begin{quickrevision}[Chapter 5 Summary]
\begin{itemize}[leftmargin=1.2em]
    \item Disaster Risk is modeled as $\text{Risk} = (\text{Hazard} \times \text{Exposure} \times \text{Vulnerability}) / \text{Capacity}$.
    \item Radiated seismic energy scales with Richter magnitude as $\log_{10} E = 4.8 + 1.5 M$ ($\approx 31.62\times$ per unit).
    \item Tropical cyclones are classified by wind speeds; super cyclonic storms exceed $222\text{ km/h}$.
    \item The Bhopal Gas Tragedy (1984) involved the release of $42\text{ tonnes}$ of toxic Methyl Isocyanate (MIC).
    \item The DM Act 2005 established a 3-tier structure: NDMA (Prime Minister), SDMA (Chief Minister), DDMA (Collector).
\end{itemize}
\end{quickrevision}
